% !PRINTED 645
\weberpage{645}
Wir wollen einem der Radicale, etwa dem $\tau_\alpha$, ein anderes Vorzeichen geben, also die Vertauschung

\weq{13}{(\tau_\alpha,\ \varepsilon^\beta \tau_\alpha)}
machen, wo $\beta$ ein beliebiger Exponent sein kann, und $\alpha$ einer der Indices $0, 1 \ldots n-2$ ist.

In (12) hat das Radical $\tau_\alpha$ den Exponenten $r_{n+\nu-\alpha-2} \equiv r_{\nu-\alpha-1}$, und also entspricht der Vertauschung (13) die Vertauschung

\weq{14}{\left(\sqrt[n]{R_\nu},\ \varepsilon^{\beta r_{\nu-\alpha-1}} \sqrt[n]{R_\nu}\right).}
Nun ist nach der Bedeutung der Zahlen $r$ allgemein $r_{\alpha+\beta} \equiv r_\alpha r_\beta \pmod{n}$, und wenn man also die Vertauschung (13) in (11) einführt, so geht $h$ in $h + \beta r_{-\alpha-1}$ über.

1. Es ruft also die Vertauschung (13) unter den Indices von $\xi$ die cyklische Permutation
\[
(h,\ h + \beta r_{-\alpha-1})
\]
hervor.

Machen wir zweitens die cyklische Permutation

\weq{15}{(\tau_0,\ \tau_1, \ldots \tau_{n-2}),}
so entspricht diese den Vertauschungen

\weq{16}{\left(\sqrt[n]{R_\nu},\ \sqrt[n]{R_{\nu+1}}\right), \quad (K_\nu, K_{\nu+1}), \quad \nu = 0, 1, 2 \ldots n-2,}
und $\xi_h$ geht über in

\[
A + \sum^{\nu} \varepsilon^{h r_\nu} K_{\nu+1} \sqrt[n]{R_{\nu+1}} = A + \sum \varepsilon^{h r_{\nu-1}} K_\nu \sqrt[n]{R_\nu}
\]
\[
= A + \sum \varepsilon^{h g^{-1} r_\nu} K_\nu \sqrt[n]{R_\nu}.
\]

2. Demnach erleiden die Indices von $\xi$ durch die Permutation (15) die Permutation
\[
(h,\ g^{-1} h).
\]

Betrachten wir nun irgend eine rationale symmetrische (oder auch nur metacyklische) Function der $\xi$:

\[
S(\xi_0,\ \xi_1 \ldots \xi_{n-1}),
\]
so erhalten wir, wenn wir die Werthe (10) einführen, daraus eine rationale Function der Radicale

\[
\tau_0,\ \tau_1, \ldots \tau_{n-2}.
\]
Diese Function ändert sich aber nicht, wenn man einem dieser Radicale einen anderen seiner $n$ Werthe giebt, und folglich muss die Function rational von $k_0, k_1 \ldots k_{n-2}$ abhängen. Wegen
%CONT>>
